% Self-contained source. Compile twice with pdflatex; no figures or .bib needed.
% Project snapshot: 2026-10-30. Evidence paths are relative to repository root.
\documentclass[10pt,letterpaper]{article}
\usepackage[T1]{fontenc}
\usepackage[utf8]{inputenc}
\usepackage{lmodern}
\usepackage[margin=0.7in]{geometry}
\usepackage{pdflscape,longtable,booktabs,array}
\usepackage[table]{xcolor}
\usepackage{amsmath,amssymb}
\usepackage{enumitem}
\usepackage{microtype}
\usepackage{fancyhdr}
\usepackage{xurl}
\usepackage[colorlinks=true,linkcolor=blue!45!black,urlcolor=blue!45!black]{hyperref}
\definecolor{ink}{HTML}{17324D}
\definecolor{pale}{HTML}{EAF0F5}
\definecolor{muted}{HTML}{465564}
\pagestyle{fancy}
\fancyhf{}
\fancyfoot[L]{\small\color{muted}Ionization Cone Automation \quad | \quad Research status: 30 October 2026}
\fancyfoot[R]{\small\thepage}
\renewcommand{\headrulewidth}{0pt}
\renewcommand{\footrulewidth}{0.3pt}
\setlength{\parindent}{0pt}
\setlength{\parskip}{6pt}
\setlength{\emergencystretch}{2em}
\setlist[itemize]{leftmargin=1.3em,itemsep=3pt,topsep=3pt}
\setlist[enumerate]{leftmargin=1.5em,itemsep=3pt,topsep=3pt}
\newcolumntype{P}[1]{>{\raggedright\arraybackslash}p{#1}}
\newcommand{\ev}[1]{\textcolor{ink}{\textbf{[#1]}}}
\newcommand{\status}[1]{\textbf{#1.}}
\newcommand{\oiii}{[O\,\textsc{iii}]}
\newcommand{\hii}{H\,\textsc{ii}}
\newcommand{\paalpha}{Pa$\alpha$}
\newcommand{\pathref}[1]{{\small\nolinkurl{#1}}}

\begin{document}
\begin{landscape}
{\LARGE\bfseries\color{ink} Science Traceability Matrix}\par
{\large Physically motivated synthetic imaging and automated AGN ionization-cone recovery}\par
\textbf{Science objective:} Determine when projected ionization-cone morphology can be recovered from emission images despite
overlapping light, cloud structure, extinction, finite resolution, and noise.

\textbf{Scope:} Recent work emphasizes 2D JWST F187N simulations and an observed NGC\,1068 reference;
HST F502N/\oiii{} and the historical 96-channel training experiment remain distinct branches.
\textbf{Status key:} \emph{Implemented} denotes a working method or diagnostic; \emph{pilot} denotes exploratory evidence;
\emph{pending} denotes a scientific validation step. None alone means validated performance on real galaxies.

\begingroup
\small
\setlength{\tabcolsep}{3pt}
\renewcommand{\arraystretch}{1.16}
\begin{longtable}{P{0.35in}P{1.35in}P{1.60in}P{1.82in}P{1.78in}P{2.10in}}
\toprule
\rowcolor{pale}
\textbf{ID} & \textbf{Science question / objective} & \textbf{Required observables / quantities} &
\textbf{Data and analysis method} & \textbf{Evidence and current status} & \textbf{Verification / remaining requirement} \\
\midrule
\endfirsthead
\multicolumn{6}{l}{\textbf{Science traceability matrix (continued)}}\\[4pt]
\toprule
\rowcolor{pale}
\textbf{ID} & \textbf{Science question / objective} & \textbf{Required observables / quantities} &
\textbf{Data and analysis method} & \textbf{Evidence and current status} & \textbf{Verification / remaining requirement} \\
\midrule
\endhead
\midrule
\multicolumn{6}{r}{\footnotesize Continued on the next page.}\\
\endfoot
\bottomrule
\endlastfoot

S1 & Define the geometry the network is intended to recover. &
Projected mask, axis PA, half/full opening, radial extent, inclination convention, and emission extent. &
Prescribed cone geometry; independent soft gas emission; explicit distinction between illumination, projected geometry, and detectable emission. &
\status{Implemented / assumed} Nominal geometric labels exist; soft emission can cross their boundaries. \ev{E1,E2} &
Audit projection and label recoverability. Do not interpret a projected mask as the intrinsic torus opening. Define any detectable-emission target separately. \\[5pt]
\midrule

S2 & Model the emitting populations that determine band morphology. &
Line/nebular spectra versus gas state; AGN, host and \hii{} gas; stellar/nuclear continua; attenuation. &
CLOUDY grids, dedicated \hii{} calculations, and prescribed spatial components; separate line, continuum, and population maps. &
\status{Pilot} Components implemented; newer converged AGN spectral pilots exist separately from older image grids. \ev{E1,E3} &
Trace each image to its actual grid. Resolve material convergence and missing-line issues; constrain relevant amplitudes and spectra. Inclusion is not validation. \\[5pt]
\midrule

S3 & Reproduce how instruments transform intrinsic emission into an image. &
Band throughput, observed wavelengths, velocity shifts, PSF, pixel scale, and noise amplitudes. &
Photon-weighted filter integration; air/vacuum conventions; spatial convolution and sampling. F187N uses an STPSF kernel and ERR-based correlated noise. &
\status{Implemented / approximate} F502N integration tests and F187N component reconstruction checks exist. \ev{E1,E4,E5} &
Keep band-specific operators distinct. Improve observation-specific PSF, detector noise, and absolute normalization where material; empirical scaling is not a luminosity calibration. \\[5pt]
\midrule

S4 & Anchor candidate geometry to observed NGC\,1068 structure. &
Sky PA, apparent angular support, radial dependence, and sensitivity to analysis choices. &
Observed HST F502N minus scaled F547M; WCS alignment with JWST F187N; angular profiles and 56 analysis trials. &
\status{Measured candidate} 55 trials return a primary lobe; PA $22.5$--$37.5^\circ$, full width $72$--$156^\circ$. \ev{E6} &
Ranges describe method sensitivity, not confidence intervals. Improve continuum subtraction/registration; optical support does not give F187N component fractions. \\[5pt]
\midrule

S5 & Quantify competing light and spatial complexity. &
$M=B/C$ in a fixed aperture; cone light fraction; area where $B>C$; structure at stated angular scales. &
Linear post-PSF component maps; exclude baseline/noise; Gaussian residual structure at $0.1$, $0.2$, $0.4$ arcsec sigma. &
\status{Implemented} Five synthetic examples characterized; 10 diagnostic tests passed. $M=1.515$--$13.231$. \ev{E5} &
Relate descriptors to actual recovery errors. Separate amplitude from structure. No universal difficulty score or observed NGC\,1068 $M$ has been established. \\
\newpage

S6 & Determine how spatial scale and noise limit recovery. &
Projected extent/PSF FWHM; pixel sampling; per-pixel cone signal/noise; angular and physical lengths. &
Optical-kernel core-width measurement; median cone/ERR; optional size conversion with a sourced angular-diameter distance. &
\status{Implemented diagnostic} Resolution and noise summaries are saved; physical distance is not inferred automatically. \ev{E5} &
Test controlled resolution/noise changes. Optical PSF width omits detector/drizzle effects; correlated pixels do not yield an independent-pixel integrated significance. \\[5pt]
\midrule

S7 & Constrain gas spectra and continuum against real data. &
Line fluxes/ratios, continuum constraints, aperture/seeing, spectral quality and atmospheric transmission. &
Published near-IR line constraints; archived ESO X-shooter spectra; preliminary telluric correction; matched slit/seeing extraction. &
\status{Pilot / discrepancies retained} Converged spectral trials and slit comparisons exist; no complete F187N decomposition is established. \ev{E3,E7} &
Use only supported wavelengths and matched apertures. Resolve line/continuum mismatches without filling gaps or arbitrary line rescaling; quantify remaining systematics. \\[5pt]
\midrule

S8 & Recover projected cone masks and angles automatically. &
Dice/IoU, PA error, opening-angle error, and failure rate versus observability and complexity. &
2D U-Net infrastructure, legacy real-image predictions, and a separate historical 96-channel cube training experiment. &
\status{Demonstrated pipeline} A saved one-epoch cube run exists; it is not an F187N training result. \ev{E8,E9} &
Prepare the current-band dataset and grouped splits; compare to a geometric baseline. Set justified performance targets before interpreting test results. \\[5pt]
\midrule

S9 & Establish generalization beyond the calibration example. &
Independent physical realizations/conditions and galaxies; domain-shift and faint-boundary performance. &
Split by underlying realization; keep filter/noise variants together. Reserve observations not used to tune generation. &
\status{Pending scientific validation} Four F187N population scenarios are exploratory; NGC\,1068 is a calibration reference. \ev{E2,E8} &
Document sampling distributions and independent holdouts. Seed variation alone on one grid/galaxy does not establish physical or observational generalization. \\[5pt]
\midrule

S10 & Make the scientific conclusions reproducible and auditable. &
Inputs, code/grid versions, parameter choices, source hashes, masks, outputs, and failure records. &
Saved configurations, component maps, manifests, signed residuals, sensitivity trials, and fresh output directories for reruns. &
\status{Implemented practice} Versioned reports and an instruction record preserve assumptions and limitations. \ev{E5,E6,E7,E10} &
Verify provenance for each release and keep failures. Evidence must distinguish a numerical check, an empirical assumption, and independent scientific validation. \\
\end{longtable}
\endgroup

\textbf{Acceptance criteria:} Numerical invariants are tested where defined. No project-wide accuracy threshold,
validated difficulty cutoff, or population-level scientific performance requirement has yet been adopted.
Evidence identifiers are resolved in the repository evidence register at the end of this document.
\end{landscape}

\section*{Summary of work completed}
\textit{The first-person summary below describes work evidenced by the repository as of 30 October 2026.
It separates implemented methods and pilot results from scientific validation still required.}

\subsection*{1. Established the simulation and learning framework}
I developed a framework for generating synthetic ionization-cone images with known projected geometry,
preparing segmentation inputs and labels, and applying U-Net models. The project includes optical
F502N/\oiii{} work, a historical spectral-cube experiment, and the more recent F187N image branch.
I preserved earlier NGC\,1068 extraction and prediction results rather than treating them as independent
ground truth. A saved 96-channel experiment completed one training epoch using six examples split into
four training, one validation, and one test example. This demonstrates the learning pipeline; it does
not demonstrate reliable recovery on real galaxies or constitute training of the current F187N model.
\ev{E8,E9}

I clarified that the generator currently \emph{prescribes} the opening angle. Emission spectra and filter
response determine how the chosen geometry appears. I distinguished intrinsic illumination geometry,
projected geometric labels, and the threshold-dependent extent visible in an emission image. This
provides an explicit interpretation for what the network is being asked to recover. \ev{E1,E10}

\subsection*{2. Added emission physics and observing effects}
I expanded the F187N pilots to include AGN-photoionized cone gas, low-ionization host gas, dedicated
\hii{} emission, nebular continuum, stellar light, and compact nuclear continuum. The images include
prescribed dust attenuation, cloud structure, velocity-dependent filter weighting, spatial blurring,
pixel sampling, and approximate correlated noise. Separate component maps are saved so that image
formation can be checked without confusing display transformations with physical quantities. \ev{E1,E4}

I also completed small converged CLOUDY spectral pilots and comparisons with published near-IR line
measurements. Those newer calculations remain distinct from the older cached grids used in existing
images. Spectral disagreements, including iron and coronal-line mismatches, were recorded rather than
removed through arbitrary line scaling. Archived ground-based spectra, preliminary atmospheric
corrections, and slit/seeing-aware model extractions provide additional observational checks, while
retaining unsupported wavelength intervals and registration uncertainties. \ev{E3,E7}

\subsection*{3. Made image ``messiness'' measurable}
I implemented separate measurements of competing light, spatial structure, resolution, and noise,
instead of assigning an unvalidated combined difficulty score. Let $C$ denote cone-associated gas,
$B$ the other explicitly listed model components, and $A$ the measurement aperture. The principal
light ratios are
\[
 M=\frac{\sum_A B\,\Delta\Omega}{\sum_A C\,\Delta\Omega},
 \qquad
 f_{\rm cone}=\frac{\sum_A C\,\Delta\Omega}{\sum_A(C+B)\,\Delta\Omega}.
\]
The maps are linear, noiseless, and already passed through the filter and PSF. The initial ratios
exclude the empirical baseline; other light is not necessarily non-AGN light. I also measure the
fraction of aperture area where $B>C$ and the normalized smoothing residual
\[
 P_s(I)=\frac{\sum_A |I-G_s*I|\,\Delta\Omega}{\sum_A I\,\Delta\Omega},
 \qquad s=0.1,\ 0.2,\ 0.4\ \mathrm{arcsec},
\]
separately for $C$ and $B$. Here $s$ is the Gaussian sigma. This structure statistic can respond to
gradients and edges as well as knots; it is not a pure cloud-clumpiness measurement. \ev{E5}

\clearpage
\subsection*{4. Characterized the existing synthetic examples}
I measured the original soft-gas pilot and four illustrative F187N population examples using the
nominal cone mask intersected with $0.3<r<3.7$ arcsec and reference-valid pixels. The measurements
below concern \emph{synthetic components}, not an observed decomposition of NGC\,1068. \ev{E5}

\begin{center}
\renewcommand{\arraystretch}{1.15}
\begin{tabular}{lrrr}
\toprule
\textbf{Synthetic example} & \textbf{$M$} & \textbf{$f_{\rm cone}$} & \textbf{Area with $B>C$}\\
\midrule
Original soft-gas pilot & 1.515 & 39.76\% & 82.53\%\\
Population reference & 5.902 & 14.49\% & 97.18\%\\
Variant 01 & 13.231 & 7.03\% & 99.65\%\\
Variant 02 & 4.221 & 19.15\% & 96.69\%\\
Variant 03 & 6.323 & 13.65\% & 96.88\%\\
\bottomrule
\end{tabular}
\end{center}

Ten tests passed for these diagnostics, covering known mixtures, normalization invariance, equal-flux
images with different structure, pixel-area weighting, undefined cases, PSF widths, and component
reconstruction. The examples differ in multiple parameters and component prescriptions, so the
table is not a controlled causal experiment or a demonstrated ranking of segmentation difficulty.

\subsection*{5. Extracted an observational morphology reference}
I used the downloaded HST F502N and F547M images to construct a signed \oiii{} emission proxy, then
measured its angular support in sky coordinates. I overlaid the candidate northern and southern
sectors on the observed JWST F187N image using FITS WCS coordinates. The extraction does not use a
U-Net prediction or a synthetic cone template to set the direction. \ev{E6}

The nominal analysis gives a northern-sector PA of $27^\circ$ east of north and an apparent full width
of $120^\circ$. Of 56 analysis trials, 55 return a primary candidate. Across those trials the PA spans
$22.5$--$37.5^\circ$ and the full width spans $72$--$156^\circ$. These are sensitivity ranges under
changes in subtraction, threshold, smoothing, radius, and center, not confidence intervals or a
measurement of the intrinsic torus opening. Three extraction tests passed, and source hashes,
signed subtraction, coverage, and WCS were checked.

The existing soft-gas configuration has PA $30^\circ$ and produces a projected full width of about
$89.44^\circ$ under its current prescription. These fall within the sensitivity ranges, but this is
not statistical validation. The observed F187N image remains total band light: I have \emph{not}
measured its cone fraction or its physical contamination ratio $M$ from this optical reference.

\subsection*{6. Current scientific position and next milestone}
The project now has a simulation framework, explicit label semantics, measurable image-complexity
descriptors, and an observational morphology reference. The next milestone is a \textbf{controlled
training pilot}, following an audit of the actual band/grid/labels and preparation of grouped data
splits. The necessary next steps are to:
\begin{enumerate}
\item Resolve emission-model issues material to the selected band and constrain component strengths
with suitable continuum or spectral data.
\item Vary contamination, structure, resolution, and noise in controlled realizations; quantify
whether the nominal geometric labels remain recoverable.
\item Compare U-Net masks and angle errors with a geometric baseline, using independent conditions
and observations for validation beyond the calibration example.
\end{enumerate}
No full 3D galaxy or torus simulation is required for that pilot. Scientific claims should remain
limited to the assumptions and validation actually tested.

\clearpage
\section*{Repository evidence register}
Paths are relative to the project root. These records support the matrix and summary; they are not
independent observations merely because they occupy different files. This document is a dated
snapshot, not a replacement for the underlying configurations and raw outputs.

\small
\begin{description}[style=nextline,leftmargin=2.5em,labelwidth=2em,itemsep=8pt]
\item[E1] \textbf{F187N physical content and label assumptions.}\\
\pathref{configs/cloudy/F187N_PHYSICS_CHECKLIST.md}\\
\pathref{configs/cloudy/ngc1068_f187n_soft_gas_v1.json}

\item[E2] \textbf{Exploratory population and empirical calibration.}\\
\pathref{configs/cloudy/F187N_REFERENCE_POPULATION.md}\\
\pathref{results/cloudy/f187n_population_pilot_v3/manifest.json}

\item[E3] \textbf{Spectral constraints and separate converged pilots.}\\
\pathref{configs/cloudy/NGC1068_SPECTRAL_CONSTRAINTS.md}\\
\pathref{results/cloudy/ngc1068_spectral_constraints_v1/fit.json}

\item[E4] \textbf{Image formation and passband verification.}\\
\pathref{configs/cloudy/NARROWBAND_TRAINING.md}\\
\pathref{scripts/cloudy/render_f187n_pilot.py}\\
\pathref{results/cloudy/f187n_psf_nrcb4/metadata.json}

\item[E5] \textbf{Complexity definitions, numerical results, and tests.}\\
\pathref{configs/cloudy/F187N_COMPLEXITY.md}\\
\pathref{results/cloudy/f187n_complexity_v1/metrics.json}\\
\pathref{tests/test_image_complexity.py}

\item[E6] \textbf{Observed cone extraction and analysis sensitivity.}\\
\pathref{configs/cloudy/NGC1068_OBSERVED_CONE.md}\\
\pathref{results/cloudy/ngc1068_observed_cone_v1/measurements.json}\\
\pathref{configs/cloudy/ngc1068_observed_cone_reference_v1.json}\\
\pathref{tests/test_observed_cone.py}

\item[E7] \textbf{Ground-based spectral and slit comparison.}\\
\pathref{configs/cloudy/NGC1068_SLIT_COMPARISON.md}\\
\pathref{results/cloudy/ngc1068_slit_comparison_v2/README.md}\\
\pathref{results/cloudy/ngc1068_telluric_dry_v1/summary.json}

\item[E8] \textbf{Historical spectral training scope and saved run.}\\
\pathref{configs/cloudy/SPECTRAL_TRAINING.md}\\
\pathref{results/cloudy/spectral_unet/run_20261008_101500_813978/history.json}\\
\pathref{results/cloudy/spectral_unet/run_20261008_101500_813978/dataset_manifest.json}

\item[E9] \textbf{Preserved earlier geometry and prediction outputs.}\\
\pathref{saved_runs/ngc1068_20260514/README.md}\\
Earlier reported angles are script/image-coordinate results, not independently verified sky PAs.

\item[E10] \textbf{Research conventions and reproducibility instructions.}\\
\pathref{AGENTS.md}
\end{description}

\end{document}
